% KHUNG BÁO CÁO – không cần sửa tệp này.
% Mỗi \phan / \gop chèn tệp tương ứng trong sections/ nếu tệp đã có,
% chưa có thì hiện dòng chữ đỏ nhắc tên tệp (xem hướng dẫn đầu main.tex).

\section{Thông tin nhóm và Phân công công việc}
\phan{01-phancong}{B}

\section{Mức độ hoàn thành}
\phan{02-hoanthanh}{B}

\section{Tìm hiểu mã nguồn xv6}

\subsection{Trường \texttt{ofile[]} của \texttt{struct proc} và \texttt{file[]} của \texttt{ftable}}
\phan{31-ofile-ftable}{A}

\subsection{Trace system call \texttt{open}}
\phan{32-open}{A}

\subsection{So sánh cách dùng file descriptor trong \texttt{open()}, \texttt{dup()} và \texttt{fork()}}
\phan{33-fd-compare}{B}

\subsection{Trace system call \texttt{fork}}
\phan{34-fork}{B}

\subsection{Các trường \texttt{pid}, \texttt{state}, \texttt{context} của \texttt{struct proc}}
\phan{35-proc-fields}{C}

\subsection{Cơ chế lập lịch: \texttt{scheduler()}, \texttt{sched()}, \texttt{yield()}}
\phan{36-scheduler}{C}

\section{Cài đặt system call \texttt{sysinfo}}

\subsection{Cấu trúc \texttt{struct sysinfo} và luồng gọi}
\phan{41-struct}{A}

\subsection{Thu thập danh sách open file}
\phan{42-files}{A}

\subsection{Thu thập danh sách process và bộ đếm \texttt{nsched}}
\phan{43-procs}{C}

\subsection{Danh sách các file đã sửa đổi}
\gop{itemize}{44-files}

\section{Lựa chọn thiết kế và giả định}
\gop{itemize}{5-thietke}

\section{Kết quả kiểm thử}
\phan{6-kiemthu}{B}

\section{Khó khăn và các cách đã thử}
\gop{dsso}{7-khokhan}

\section{Những khái niệm chưa hiểu rõ}
\gop{dsso}{8-chuahieu}

\nocite{*}
